% SIAM SIURO Article Template — Group 05 Final Report
\documentclass[final,onefignum,onetabnum]{siuro250211}

\usepackage{amsfonts}
\usepackage{graphicx}
\usepackage{epstopdf}
\usepackage{algorithmic}
\ifpdf
  \DeclareGraphicsExtensions{.eps,.pdf,.png,.jpg}
\else
  \DeclareGraphicsExtensions{.eps}
\fi

\usepackage{enumitem}
\setlist[enumerate]{leftmargin=.5in}
\setlist[itemize]{leftmargin=.5in}

\usepackage{tabularx}
\newcolumntype{Y}{>{\raggedright\arraybackslash}X}

\newcommand{\creflastconjunction}{, and~}

\newsiamremark{remark}{Remark}
\newsiamremark{hypothesis}{Hypothesis}
\crefname{hypothesis}{Hypothesis}{Hypotheses}
\newsiamthm{claim}{Claim}
\newsiamremark{fact}{Fact}
\crefname{fact}{Fact}{Facts}

\setcounter{footnote}{1}

\headers{Bullwhip Effect in Sri Lanka's Rice Supply Chain}{Group 05}

\title{An Agent-Based Model of the Bullwhip Effect in Sri Lanka's
North Central Province Rice Supply Chain During the 2022 Economic Crisis}

\author{H.W.S.N. Hathurusinghe
\and U.A.D.L.T. Subasinghe
\and D.A.A.K. Dewasingha
\and U.K.D.I. Dilshan
\and T.P.N. Fernando}

\dedication{\small\textit{Project advisor: Dr. S. C. Surasinghe}}

\usepackage{amsopn}
\DeclareMathOperator{\diag}{diag}

\ifpdf
\hypersetup{
  pdftitle={An Agent-Based Model of the Bullwhip Effect in Sri Lanka's Rice Supply Chain},
  pdfauthor={Group 05}
}
\fi

\begin{document}

\maketitle

\section{Introduction}
\label{sec:intro}

The Bullwhip Effect describes the progressive amplification of order variance as information
travels upstream through a supply chain: small, easily absorbed fluctuations in end-consumer
demand translate into increasingly large and costly order swings at the retailer, wholesaler, and
producer tiers \cite{LeePadmanabhanWhang1997}. The effect is well documented in stable-demand
industrial settings, but its behaviour under acute, compounding crisis conditions---where supply
capacity, transport, and demand are simultaneously disrupted---is comparatively underexplored.

Sri Lanka's 2022 economic crisis provides a real-world natural experiment in exactly this
setting. Within a period of months, the country experienced (i) a nationwide fertilizer import
ban that sharply reduced paddy yields, (ii) a severe fuel shortage that constrained transport and
logistics capacity across the supply chain, and (iii) a period of consumer panic buying driven by
food-security anxiety. This paper models the resulting demand and supply disruption in the rice
supply chain of Sri Lanka's North Central Province, the country's principal rice-producing
region, using an agent-based simulation calibrated to the timing and magnitude of these three
shocks.

The model's core ordering logic is based on the Order-Up-To policy and return-allowance
framework of Chatfield and Pritchard \cite{ChatfieldPritchard2013}, extended with an adaptive
Government agent that is not present in the original formulation. Three research questions
motivate the experimental design (\cref{sec:experiments}):

\begin{enumerate}
\item How do the fertilizer ban, fuel shortage, and consumer panic signal interact to amplify
demand variance across the supply chain?
\item Which government interventions most effectively reduce bullwhip amplification during a
crisis?
\item How sensitive is the magnitude and duration of the amplification cascade to model
parameter uncertainty and to the timing of policy activation?
\end{enumerate}

The remainder of this paper is organized as follows. \Cref{sec:related} situates the model
relative to prior bullwhip effect literature. \Cref{sec:model} describes the agent-based model
implementation. \Cref{sec:experiments} details the experimental design and data collection
protocol. \Cref{sec:analysis} describes the data analysis methodology. \Cref{sec:results} presents
the simulation results, and \cref{sec:discussion} discusses the findings, limitations, and
concludes.

\section{Related Work}
\label{sec:related}

The Bullwhip Effect was formalized by Lee, Padmanabhan, and Whang \cite{LeePadmanabhanWhang1997},
who identified demand-signal processing, order batching, price fluctuation, and rationing/shortage
gaming as its four principal operational causes. Chatfield and Pritchard
\cite{ChatfieldPritchard2013} extend this literature by incorporating product returns into a
multi-echelon simulation, showing that permitting a nonzero return allowance increases
order-variance amplification relative to a no-returns baseline, with the effect strongest at the
most upstream echelon; the return allowance parameter $\rho$ and the Order-Up-To replenishment
policy used in the present model are adopted directly from that work.

Dominguez, Cannella, and Framinan \cite{DominguezCannellaFraminan2015} study how supply chain
network structure---the number of echelons, number of nodes, and distribution of links---affects
bullwhip severity under both stationary and highly volatile demand regimes, using a multi-agent
simulation and a formal design-of-experiments approach; their finding that structural factors
beyond echelon count matter under demand shocks motivates this paper's attention to a fixed,
realistically tiered four-echelon structure rather than a stylized linear chain.

Moyaux, Chaib-draa, and D'Amours \cite{MoyauxChaibdraaDamours2007} examine information sharing as
a coordination mechanism for bullwhip reduction, separating raw consumer demand from
downstream ordering adjustments and showing that sharing true demand information upstream can
substantially dampen amplification; this motivates the present model's Government agent, which
can be interpreted as a coarse, centrally observed proxy for the food-security signal that
individual supply chain tiers do not otherwise observe.

Qu and Raff \cite{QuRaff2017} compare centralized and decentralized inventory control regimes and
show that centralization does not uniformly reduce the bullwhip effect; this result is directly
relevant to interpreting the present paper's finding that government intervention effectiveness
depends critically on activation timing rather than on the mere presence of centralized policy
control.

\section{Model Implementation}
\label{sec:model}

The model is implemented as a discrete-time agent-based simulation in NetLogo 6.4/7.0, with a time
step of one week and a horizon of 156 steps (three calendar years, 2021--2023). Development
followed an incremental approach: a minimal two-tier prototype (Consumer, Farmer) with a fixed
ordering rule was validated first, after which the Retailer and Wholesaler tiers, the return
allowance mechanism, the three crisis shocks, and the adaptive Government agent were added in
sequence, with the model's behaviour checked against the expected qualitative pattern at each
stage. The model file includes inline comments documenting each procedure and an Info tab with
the full parameter list, units, and calibration sources.

\subsection{Agents}

Five agent types are implemented:

\begin{itemize}
\item \textbf{Consumer}: generates weekly demand, subject to a panic-buying multiplier during the
historically calibrated panic window.
\item \textbf{Retailer}, \textbf{Wholesaler}: intermediate tiers, each applying an Order-Up-To
replenishment policy against a lead time pipeline of orders placed upstream.
\item \textbf{Farmer}: the production tier, subject to capacity shocks from the fertilizer ban and
fuel shortage, and the source of the food security index observed by the Government agent.
\item \textbf{Government}: an adaptive policy agent that monitors the food security index and, once
it falls below a threshold $\Phi^{*}$, activates an order-amplification cap on downstream
ordering behaviour.
\end{itemize}

\subsection{Ordering Policy}

Each tier orders according to an Order-Up-To policy, following Chatfield and Pritchard
\cite{ChatfieldPritchard2013}. Let $D_t$ denote observed demand at time $t$, $\hat{D}_t$ the
exponentially smoothed demand forecast with smoothing parameter $\alpha$, and $L$ the tier's lead
time (in weeks). The order-up-to level $S_t$ and base order quantity $O_t^{\text{base}}$ are
computed as

\begin{subequations}
\begin{align}
\hat{D}_t &= \alpha D_t + (1-\alpha)\hat{D}_{t-1}, \label{eq:forecast}\\
S_t &= \hat{D}_t \cdot (L + 1) + z\,\sigma_{D}\sqrt{L}, \label{eq:orderupto}\\
O_t^{\text{base}} &= \max\big(0,\; S_t - I_t - \textstyle\sum_{k} P_{t,k}\big), \label{eq:orderqty}
\end{align}
\end{subequations}

where $I_t$ is on-hand inventory, $P_{t,k}$ is the pipeline of outstanding orders at lag $k$, and
$z\sigma_D\sqrt{L}$ is a safety-stock term. Lead times are tiered across the chain (2, 3, and 4
weeks for the retailer, wholesaler, and farmer tiers respectively), consistent with increasing
logistics complexity moving upstream.

\subsection{Returns}

Following \cite{ChatfieldPritchard2013}, a return allowance $\rho$ introduces an unpredictable
additional source of order-quantity volatility, rather than a predictable adjustment to on-hand
inventory. The final order quantity is
\begin{equation}
O_t = \max\big(0,\; O_t^{\text{base}} + \rho \cdot O_t^{\text{base}} \cdot \varepsilon_t\big),
\qquad \varepsilon_t \sim \mathcal{N}(0,1),
\label{eq:returns}
\end{equation}
so that at $\rho=0$ the model reduces exactly to the no-returns case, and increasing $\rho$ scales
the magnitude of an unpredictable shock applied directly to the order signal, consistent with
returns functioning as an additional, disruptive information-distortion mechanism rather than a
smoothing buffer.

\subsection{Crisis Shocks}

Three shocks are applied at historically calibrated times (\cref{tab:shocks}), based on Central
Bank of Sri Lanka reporting \cite{CBSL2022}, Department of Census and Statistics paddy yield data
\cite{DCS2022Maha}, and World Food Programme price monitoring \cite{WFP2022}.

\begin{table}[htbp]
\footnotesize
\caption{Crisis shock parameters and calibration sources.}
\label{tab:shocks}
\begin{center}
\begin{tabularx}{\textwidth}{|l|c|c|Y|}
\hline
\textbf{Shock} & \textbf{Timing (week)} & \textbf{Multiplier} & \textbf{Source} \\
\hline
Fertilizer import ban & 17 & 0.627 & DCS Maha yield decline (4{,}492$\to$2{,}819 kg/ha, $\approx$37\%) \\
\hline
Fuel shortage & 65 & 0.75 & Estimated logistics capacity reduction \\
\hline
Consumer panic buying & 74--90 & 1.02/week & Calibrated to WFP price-spike timing \\
\hline
\end{tabularx}
\end{center}
\end{table}

The fertilizer-shock multiplier reflects a scientifically grounded, data-derived value; the fuel
shortage and panic-buying multipliers are acknowledged estimates in the absence of comparably
precise administrative data, and are treated as such in the limitations discussion
(\cref{sec:discussion}).

\subsection{Government Intervention}

The food security index $\Phi_t$ is computed as a weighted combination of farmer production
capacity and farmer inventory,
\begin{equation}
\Phi_t = 0.5 \left(\frac{C_t}{100}\right) + 0.5\left(\frac{I^{F}_t}{200}\right),
\label{eq:fsi}
\end{equation}
where $C_t$ is farmer production capacity (as a percentage of nominal) and $I^{F}_t$ is farmer
inventory. The Government agent activates an order-amplification cap once $\Phi_t$ falls below a
threshold $\Phi^{*}$ (default 0.40, following WFP food-security reporting convention
\cite{WFP2022}).

\subsection{Outcome Metric}

The primary outcome metric is the Bullwhip Ratio at a given tier, defined as the ratio of that
tier's order-quantity variance to consumer demand variance over the run,
\begin{equation}
\mathrm{BR}_{\text{tier}} = \frac{\operatorname{Var}(O^{\text{tier}})}{\operatorname{Var}(D)}.
\label{eq:bwr}
\end{equation}
The farmer-tier Bullwhip Ratio, being furthest upstream, is the primary metric reported in this
paper, consistent with the theoretical expectation that amplification is largest at the tier
farthest from the consumer \cite{LeePadmanabhanWhang1997}.

\section{Experimental Design and Data Collection}
\label{sec:experiments}

Four BehaviorSpace experiments were implemented directly in the NetLogo model file, in addition
to one supplementary threshold-sensitivity experiment. All experiments use the 156-step time
horizon and record output in NetLogo's table export format, giving one row per (run, time step)
pair. \Cref{tab:experiments} summarizes the design.

\begin{table}[htbp]
\footnotesize
\caption{Experimental design summary.}
\label{tab:experiments}
\begin{center}
\begin{tabularx}{\textwidth}{|l|Y|c|c|l|}
\hline
\textbf{Exp.} & \textbf{Design} & \textbf{Scen.} & \textbf{Runs} & \textbf{RQ} \\
\hline
E0 & Baseline validation, all shocks off, government off & 1 & 10 & Val. \\
\hline
E1 & Shock decomposition, $2\times2\times2$ factorial (fertilizer $\times$ fuel $\times$ panic) & 8 & 400 & RQ1 \\
\hline
E2 & Intervention $\times$ returns $\times$ panic severity, $2\times3\times3$ factorial & 18 & 900 & RQ2, RQ3 \\
\hline
E3 & One-factor-at-a-time sensitivity ($\alpha$, $\rho$, panic-rate, $\Phi^{*}$, $\pm20\%$) & 12 & 360 & RQ3 \\
\hline
FSI check & $2\times3$ factorial (government $\times$ $\Phi^{*}\in\{0.2,0.4,0.6\}$) & 6 & 300 & RQ2, RQ3 \\
\hline
\multicolumn{3}{|r|}{\textbf{Total}} & \textbf{1{,}970} & \\
\hline
\end{tabularx}
\end{center}
\end{table}

E0 confirms that, absent any crisis shock, the model converges to a stable, low-variance
order-up-to equilibrium (mean farmer-tier Bullwhip Ratio 51.5, $\sigma=6.5$, $n=10$), establishing
a validity baseline against which E1--E3 are interpreted. E1 isolates the individual and joint
contribution of each of the three shocks to demand-variance amplification (RQ1). E2 is the primary
experiment for comparing farmer-tier Bullwhip Ratios with and without government intervention, for
testing whether permitting returns amplifies variance as reported by
\cite{ChatfieldPritchard2013}, and for assessing intervention robustness under more severe demand
shocks (RQ2, RQ3). E3 evaluates robustness to parameter uncertainty by varying each of four
parameters independently by $\pm20\%$ around its calibrated default (RQ3). The supplementary
FSI-threshold experiment tests whether the activation threshold $\Phi^{*}$ itself, rather than
merely the presence of the government switch, determines intervention effectiveness.

Fifty replications per scenario were used for E1 and E2, consistent with the replication protocol
adopted from \cite{ChatfieldPritchard2013}; thirty replications were used for the E3 sensitivity
sweeps; fifty replications per scenario were used for the FSI-threshold check.

\section{Data Analysis and Visualization}
\label{sec:analysis}

Simulation output was exported from NetLogo's BehaviorSpace in table format (one row per run/time
step pair) and analyzed in Python using \texttt{pandas} for data cleaning and aggregation,
\texttt{scipy.stats} for significance testing, and \texttt{matplotlib} for visualization. Data
cleaning consisted of removing BehaviorSpace's metadata header rows, standardizing column names,
and casting parameter and metric columns to their correct types.

Group differences were tested using Welch's $t$-test (unequal variances), selected because
government-on runs exhibit substantially lower variance than government-off runs at effective
activation thresholds, violating the equal-variance assumption of a standard Student's $t$-test.
Four visualization types were produced: time-series plots of the farmer-tier Bullwhip Ratio
faceted by threshold, a boxplot of final-step ratio distributions by scenario, a heatmap of mean
ratios across the full government$\times$threshold grid, and a bar chart of the E1 shock-
decomposition means with standard-error bars.

\section{Results}
\label{sec:results}

Results below draw on four BehaviorSpace experiments. E0 (10 runs) validates baseline model
stability absent any shock. E1 (400 runs) decomposes the individual and joint contribution of the
three crisis shocks with government intervention off. The primary intervention experiment (E2,
900 runs) compares government-on and government-off replications across a $3\times3$ grid of
return allowance $\rho\in\{0,0.1,0.2\}$ and panic-rate severity, at the calibrated default
activation threshold ($\Phi^{*}=0.40$). A supplementary threshold-sensitivity experiment (300
runs) extends the government comparison across three activation thresholds
($\Phi^{*}\in\{0.2,0.4,0.6\}$). All experiments report the farmer-tier Bullwhip Ratio
(\cref{eq:bwr}) at the final simulation step unless otherwise noted.

\subsection{Shock Decomposition (E1)}

Holding government intervention off, no individual shock produced a statistically significant
main effect on the final-step farmer-tier Bullwhip Ratio when averaged across the other two
factors: fertilizer ban ($p=0.88$), fuel shortage ($p=0.057$), consumer panic ($p=0.67$)
(\cref{fig:e1}). Direct inspection of per-tick consumer demand confirmed that the panic-rate
multiplier (a 2\% mean shift) is applied correctly by the model, but its magnitude is small
relative to the model's intrinsic demand volatility ($\sigma=10$ on a mean of 50), rendering its
contribution to the cumulative variance ratio statistically undetectable at this calibration. The
fuel shortage, a 25\% capacity reduction, comes closest to significance, consistent with its
larger direct effect on farmer production capacity.

\begin{figure}[htbp]
  \centering
  \includegraphics[width=0.95\textwidth]{e1_shock_decomposition}
  \caption{Mean farmer-tier Bullwhip Ratio across all eight shock combinations in E1 (government
  off), with standard-error bars. No individual shock shows a statistically significant main
  effect (all $p>0.05$).}
  \label{fig:e1}
\end{figure}

\subsection{Return Allowance Effect (E2)}

With government intervention off, increasing the return allowance $\rho$ increases the mean
farmer-tier Bullwhip Ratio (\cref{fig:rho}): 72.8 ($\rho=0$), 74.0 ($\rho=0.1$, not significant
relative to $\rho=0$, $p=0.47$), and 86.5 ($\rho=0.2$, $p<0.0001$ relative to $\rho=0$). This
confirms, in direction, the central finding of Chatfield and Pritchard \cite{ChatfieldPritchard2013}
that permitting returns increases bullwhip amplification, though the effect in this model is only
clearly detectable at the higher return allowance level tested.

\begin{figure}[htbp]
  \centering
  \includegraphics[width=0.7\textwidth]{rho_effect}
  \caption{Mean farmer-tier Bullwhip Ratio as a function of return allowance $\rho$ (government
  off). Error bars show standard error of the mean ($n=150$ per level).}
  \label{fig:rho}
\end{figure}

\subsection{Primary Intervention Effect (E2)}

At the calibrated default settings ($\rho=0.1$, panic-rate$=1.02$), government intervention
reduces the mean farmer-tier Bullwhip Ratio by 61.3\%, from 74.53 (government off, $n=50$) to
28.85 (government on, $n=50$; Welch's $t$-test, $p=9.9\times10^{-29}$). This effect is consistent
across all nine combinations of $\rho$ and panic-rate severity tested in the full E2 factorial
design.

\subsection{Threshold Sensitivity}

The supplementary FSI-threshold experiment corroborates the primary intervention effect and shows
it is threshold-dependent (\cref{tab:threshold}, \cref{fig:threshold,fig:heatmap,fig:boxplot}).

\begin{table}[htbp]
\footnotesize
\caption{Farmer-tier Bullwhip Ratio by FSI activation threshold, government on vs.\ off (50 runs
per cell). $p$-values from Welch's $t$-test (unequal variances).}
\label{tab:threshold}
\begin{center}
\begin{tabular}{|c|c|c|c|c|c|}
\hline
\textbf{$\Phi^{*}$} & \textbf{Off (mean)} & \textbf{On (mean)} & \textbf{Reduction} & \textbf{$p$-value} & \textbf{Significant?} \\
\hline
0.2 & 77.07 & 74.45 & 3.4\% & 0.393 & No \\
\hline
0.4 & 79.26 & 27.52 & 65.3\% & $9.2\times10^{-25}$ & Yes \\
\hline
0.6 & 72.52 & 23.93 & 67.0\% & $8.5\times10^{-31}$ & Yes \\
\hline
\end{tabular}
\end{center}
\end{table}

At $\Phi^{*}=0.2$, intervention triggers only once the food security index has nearly collapsed,
too late to prevent order amplification, and the difference from the no-intervention baseline is
not statistically significant. At $\Phi^{*}=0.4$ and $0.6$, intervention activates earlier in the
crisis timeline and substantially suppresses the farmer-tier Bullwhip Ratio, cutting it by
roughly 65--67\%.

\begin{figure}[htbp]
  \centering
  \includegraphics[width=0.95\textwidth]{bullwhip_by_threshold}
  \caption{Farmer-tier Bullwhip Ratio over time, government on vs.\ off, faceted by FSI
  activation threshold. Dotted vertical lines mark the fertilizer-ban, fuel-shortage, and
  panic-buying shock onsets.}
  \label{fig:threshold}
\end{figure}

\begin{figure}[htbp]
  \centering
  \includegraphics[width=0.75\textwidth]{heatmap_mean_bwr}
  \caption{Mean farmer-tier Bullwhip Ratio, government intervention status $\times$ FSI activation
  threshold.}
  \label{fig:heatmap}
\end{figure}

\begin{figure}[htbp]
  \centering
  \includegraphics[width=0.85\textwidth]{boxplot_final_distribution}
  \caption{Final-step Bullwhip Ratio distribution by scenario (50 runs per box), government on
  vs.\ off, across all three FSI activation thresholds.}
  \label{fig:boxplot}
\end{figure}

\section{Discussion and Conclusion}
\label{sec:discussion}

\subsection{Summary of Main Findings}

This study modeled the Bullwhip Effect in Sri Lanka's North Central Province rice supply chain
during the 2022 economic crisis, testing whether an adaptive, food-security-triggered government
intervention reduces upstream order-variance amplification, and whether permitting returns
increases it. Both hypotheses are supported. Government intervention reduces the farmer-tier
Bullwhip Ratio by 61.3\% at calibrated default settings, though this benefit is conditional on
early activation: intervention is statistically ineffective when triggered too late in the crisis
timeline ($\Phi^{*}=0.2$). The return allowance $\rho$ increases the farmer-tier Bullwhip Ratio,
consistent in direction with Chatfield and Pritchard \cite{ChatfieldPritchard2013}, though the
effect is only clearly detectable at the higher return level tested. The shock-decomposition
experiment (E1) found that none of the three individual crisis shocks produced a statistically
significant effect in isolation, suggesting the shocks' contribution to amplification is better
understood as arising from their interaction with the ordering policy and government-intervention
mechanism (E2) than from their individual, isolated presence.

\subsection{Interpretation and Implications}

The mechanism behind the observed amplification is consistent with standard bullwhip theory
\cite{LeePadmanabhanWhang1997}: each tier's Order-Up-To policy reacts to smoothed local demand
rather than true consumer demand, so successive tiers over- and under-correct, with lead time and
the return allowance $\rho$ compounding the effect moving upstream. The Government agent
interrupts this compounding process by capping order amplification once the food security index
signals crisis conditions---functioning, in effect, as a coarse form of the demand-information
sharing mechanism studied by Moyaux, Chaib-draa, and D'Amours \cite{MoyauxChaibdraaDamours2007}.
The threshold-dependence result indicates that this substitute signal only helps if it fires while
farmer capacity and inventory are still recoverable; once the index has nearly collapsed, the cap
can no longer prevent the amplification that has already propagated through the chain.

The E1 finding that no individual shock reaches significance in isolation is itself informative:
it suggests that in this model, crisis-driven bullwhip amplification is not simply the sum of
independent shock contributions, but emerges from how the ordering policy, lead-time structure,
and (where active) government response interact with the shocks in combination---a point directly
relevant to RQ1 and worth further decomposition in future work (\cref{sec:future}).

\subsection{Comparison to Previous Work}

Chatfield and Pritchard \cite{ChatfieldPritchard2013} find that returns increase bullwhip
amplification most strongly at the most upstream echelon; the present model's farmer tier
(analogous to their Factory echelon) shows the same qualitative direction, though the magnitude of
increase detected here is more modest and only reaches significance at $\rho=0.2$ rather than
scaling smoothly across all tested levels. Dominguez, Cannella, and Framinan
\cite{DominguezCannellaFraminan2015} find that supply chain structural factors beyond echelon
count materially affect bullwhip severity under demand shocks; the present four-tier, tiered-
lead-time structure produces amplification magnitudes broadly consistent with their volatile-
demand scenarios, though a direct structural comparison was outside this paper's scope. Qu and
Raff \cite{QuRaff2017} find that centralized inventory control does not uniformly outperform
decentralized control; the present paper's threshold-dependence finding is a specific instance of
this general result, showing that a nominally centralized intervention can fail entirely under
poor parameterization.

\subsection{Limitations}

Several limitations should be noted. First, the implemented government agent uses a single
adaptive policy lever (an order-amplification cap), rather than the three discrete levers (price
controls, import relief, rationing) originally proposed. Second, two of the three shock
multipliers (fuel shortage, consumer panic rate) are calibrated to estimated rather than directly
measured values, in the absence of comparably precise administrative data to the DCS paddy yield
series used for the fertilizer shock. Third, the return allowance effect (\cref{fig:rho}) was only
clearly significant at $\rho=0.2$; the step from $\rho=0$ to $\rho=0.1$ did not reach significance,
so the relationship should not be read as strictly linear across the tested range. Fourth, the E3
one-factor-at-a-time sensitivity sweep for $\rho$ specifically was executed prior to a
subsequent correction to the returns mechanism (\cref{eq:returns}) and its results are therefore
not reported here; the $\alpha$, panic-rate, and $\Phi^{*}$ sweeps in E3 do not depend on the
returns mechanism and remain valid, but a repeated $\rho$-sensitivity sweep under the corrected
mechanism is identified as outstanding work. Fifth, a supplementary AnyLogic Personal Learning
Edition model was also built for cross-checking purposes; it cannot replicate NetLogo's full
50-run BehaviorSpace design due to platform replication limits. All statistical claims and all
figures in this report are drawn exclusively from NetLogo BehaviorSpace output, analyzed and
visualized in Python (\cref{sec:analysis}).

\subsection{Future Directions}
\label{sec:future}

Three extensions follow naturally from these limitations. First, re-running the E3 $\rho$-
sensitivity sweep under the corrected returns mechanism would complete the parameter-robustness
picture for RQ3. Second, replacing the single order-amplification cap with the originally proposed
three-lever policy space (price controls, import relief, rationing) would allow testing whether
more targeted interventions outperform the current coarse cap. Third, given that no individual
shock in E1 reached significance in isolation, a follow-up factorial analysis explicitly testing
two- and three-way shock interaction terms (rather than main effects alone) would help clarify
whether amplification in this model is driven by shock interactions, by interaction with the
ordering policy structure, or by the government-intervention mechanism's response to combined
shocks.

\subsection{Conclusion}

This paper models the Bullwhip Effect in Sri Lanka's North Central Province rice supply chain
during the 2022 economic crisis using an agent-based simulation grounded in the Order-Up-To and
returns framework of Chatfield and Pritchard \cite{ChatfieldPritchard2013}. An adaptive government
agent, activating an order-amplification cap once a food security index crosses a threshold,
reduces farmer-tier order-variance amplification by roughly 61--67\% when activated early enough
in the crisis timeline, but provides no statistically detectable benefit when activated too late.
Permitting returns increases farmer-tier amplification, consistent in direction with the cited
literature. These findings suggest that the design of food-security policy triggers---specifically,
their activation timing relative to a compounding supply shock---is at least as consequential as
the decision to intervene at all.

\section{Individual Contribution Statement}
\label{sec:contributions}

\begin{itemize}[leftmargin=.3in]

\item[] \textbf{U.K.D.I. Dilshan, 16939}
\begin{itemize}
\item \textit{Submodel implemented or analyzed:} Base ordering model, following Chatfield and
Pritchard \cite{ChatfieldPritchard2013}. Implemented the Order-Up-To replenishment policy
(\cref{eq:forecast,eq:orderupto,eq:orderqty}), the tiered lead time pipelines, and the return
allowance shock mechanism (\cref{eq:returns}) in NetLogo, including the diagnosis and correction of
two implementation bugs (a consumer-demand overwrite affecting the panic-rate shock, and an
inventory-buffering returns mechanism producing the wrong effect direction relative to
\cite{ChatfieldPritchard2013}).
\item \textit{Analysis or visualization contribution:} Led NetLogo BehaviorSpace configuration and
execution for experiments E0--E3 and the FSI-threshold check.
\item \textit{Writing contribution:} Wrote \cref{sec:intro,sec:model,sec:experiments}.
\item \textit{Understanding of project:} This project models how the 2022 economic crisis in Sri
Lanka disrupted the rice supply chain in the North Central Province, using an agent-based
simulation of five interacting agents: Consumer, Retailer, Wholesaler, Farmer, and an adaptive
Government agent. My contribution was the ordering logic that every non-consumer agent uses each
week: each tier forecasts demand, compares it against its target inventory level, and places an
order to close the gap, with a delay of several weeks before that order arrives (the lead time).
I also implemented the mechanism by which returns of unsold or excess stock introduce extra,
unpredictable swings into that ordering decision, following Chatfield and Pritchard's finding that
allowing returns increases the bullwhip effect rather than dampening it. During development I
found and corrected two bugs: one where the panic-buying shock was silently discarded before it
could affect demand, and one where returns were implemented in a way that made the model's results
run backwards relative to the cited paper. Fixing both required not just changing code but
re-verifying, with real BehaviorSpace output, that the fix produced the expected direction of
effect before trusting any downstream result. This project taught me that an agent-based model is
only as trustworthy as its weakest mechanism, and that a plausible-looking result can still be
wrong for structural reasons that only show up when you check it against the theory it's supposed
to reproduce.
\end{itemize}

\item[] \textbf{U.A.D.L.T. Subasinghe, 16652}
\begin{itemize}
\item \textit{Submodel implemented or analyzed:} Reviewed Dominguez, Cannella, and Framinan
\cite{DominguezCannellaFraminan2015} on supply chain structural factors; contributed to the design
and justification of the model's four-tier chain structure.
\item \textit{Analysis or visualization contribution:} Cross-checked the model's four-tier,
tiered-lead-time chain structure against the structural factors identified in
\cite{DominguezCannellaFraminan2015} (echelon count, lead-time distribution across tiers) to
confirm the model's design choices were consistent with that paper's findings on which structural
factors most affect bullwhip severity under volatile demand.
\item \textit{Writing contribution:} Contributed to \cref{sec:related}.
\item \textit{Understanding of project:} My paper, Dominguez, Cannella, and Framinan (2015),
studies how a supply chain's structure---not just how many tiers it has, but how those tiers are
connected and how lead times are distributed among them---affects how badly a demand shock gets
amplified as it travels upstream. This directly shaped a design decision in our own model: we
built a fixed four-tier chain (Consumer, Retailer, Wholesaler, Farmer) with different lead times
at each tier (2, 3, and 4 weeks respectively) rather than a simplified chain with uniform lead
times, because the literature shows that structural detail like this changes the outcome under
crisis conditions, not just under normal demand. My role was to make sure this part of our model's
design was defensible against the literature, and to help explain, in the Related Work section,
why a stylized linear chain would not have been an adequate choice for a project specifically
about a crisis scenario.
\end{itemize}

\item[] \textbf{D.A.A.K. Dewasingha, 16935}
\begin{itemize}
\item \textit{Submodel implemented or analyzed:} FSI-threshold sensitivity experiment design and
execution.
\item \textit{Analysis or visualization contribution:} Conducted the Python-based statistical
analysis (\texttt{pandas}, \texttt{scipy}, \texttt{matplotlib}) of BehaviorSpace CSV exports,
including Welch's $t$-test significance testing, and produced the visualizations in
\cref{fig:e1,fig:rho,fig:threshold,fig:heatmap,fig:boxplot}.
\item \textit{Writing contribution:} Wrote \cref{sec:analysis} and the Results subsections in
\cref{sec:results}.
\item \textit{Understanding of project:} My contribution focused on turning the raw NetLogo
BehaviorSpace output into the statistical evidence the report relies on. Each experiment produces
one CSV row per simulation run per week, so the first step was always cleaning that data (removing
BehaviorSpace's header rows, fixing column types) before any analysis could happen. I designed the
supplementary FSI-threshold experiment specifically to test whether the government agent's
activation threshold, not just whether the government switch was on or off, determined whether
intervention actually helped, and used Welch's $t$-test rather than a standard $t$-test because the
government-on and government-off groups have very different variances, which would otherwise bias
the significance test. The most important finding I helped establish was that government
intervention only works if it activates early enough in the crisis timeline: at a threshold of
0.2, the food security index has already nearly collapsed by the time the government acts, and the
result is statistically indistinguishable from doing nothing at all. That finding shaped how the
whole team framed the project's conclusion, from "government intervention helps" to the more
precise and more useful claim that intervention timing is what actually matters.
\end{itemize}

\item[] \textbf{T.P.N. Fernando, 16941}
\begin{itemize}
\item \textit{Submodel implemented or analyzed:} Reviewed Qu and Raff \cite{QuRaff2017} on
centralized versus decentralized inventory control.
\item \textit{Analysis or visualization contribution:} Compared the model's government-intervention
mechanism against the centralized-control framework in \cite{QuRaff2017}, helping interpret why
the FSI-threshold experiment's finding (intervention effectiveness depends on activation timing,
not merely on the presence of a centralized policy) is consistent with that paper's broader result
that centralization does not uniformly outperform decentralized control.
\item \textit{Writing contribution:} Contributed to the Comparison to Previous Work and
policy-implications discussion in \cref{sec:discussion}.
\item \textit{Understanding of project:} My paper, Qu and Raff (2017), compares centralized and
decentralized inventory control across a supply chain and finds that centralizing control does not
automatically reduce the bullwhip effect; it depends on how and when that central control is
exercised. This connects directly to our project's government agent, which is a form of
centralized intervention: it observes a food-security signal across the whole chain and caps
orders once conditions look bad enough. Our own results echo Qu and Raff's point almost exactly:
the government agent only reduces the bullwhip effect when it activates early, at a threshold of
0.4 or 0.6; at a more conservative threshold of 0.2, it provides no statistically measurable
benefit at all, even though it is still, nominally, a centralized intervention. My contribution was
connecting this literature to that specific result, so the report doesn't just report "government
intervention works" as a flat finding, but explains why a poorly timed centralized policy can fail
even when a well-timed one succeeds, which is the more accurate and more useful conclusion.
\end{itemize}

\item[] \textbf{H.W.S.N. Hathurusinghe, 16942}
\begin{itemize}
\item \textit{Submodel implemented or analyzed:} Reviewed Moyaux, Chaib-draa, and D'Amours
\cite{MoyauxChaibdraaDamours2007} on information sharing as a coordination mechanism; contributed
to the conceptual framing of the Government agent.
\item \textit{Analysis or visualization contribution:} Evaluated the Government agent's food
security index mechanism as an implementation of the information-sharing coordination principle
described in \cite{MoyauxChaibdraaDamours2007}, comparing it against that paper's finding that
sharing true demand information upstream dampens bullwhip amplification.
\item \textit{Writing contribution:} Contributed to \cref{sec:related} and the Government
Intervention subsection of \cref{sec:model}.
\item \textit{Understanding of project:} My paper, Moyaux, Chaib-draa, and D'Amours (2007), studies
how sharing information about true consumer demand, rather than letting each tier only see the
orders placed by its immediate neighbor, can reduce the bullwhip effect. Individual tiers in a
supply chain normally only see a distorted version of demand, filtered through every tier below
them; sharing the real signal upstream removes some of that distortion. This paper directly shaped
how we designed the Government agent in our model. Rather than letting the government see raw
orders like every other tier does, we gave it its own food security index, computed from the
farmer's actual production capacity and inventory, so it has a clearer, more direct signal of the
crisis than any single supply chain tier would have on its own. My contribution was making sure
this design choice was grounded in the coordination-mechanism literature rather than just being an
arbitrary rule, and helping write up why an information-based intervention like this is a
reasonable stand-in for the more detailed three-lever policy (price controls, import relief,
rationing) that was originally proposed but not fully implemented.
\end{itemize}

\end{itemize}

All members reviewed and approved the final report and jointly prepared the GitHub repository,
\href{https://github.com/sara316129/srilanka-rice-supply-chain-bullwhip}{\textit{Bullwhip Effect
in Sri Lanka's Rice Supply Chain}}, containing the model, experiment CSV outputs, analysis
scripts, and this document.

\section{Data and Code Availability}
\label{sec:data-availability}

The NetLogo model, BehaviorSpace experiment definitions, exported CSV data, and Python analysis
scripts are available in the project repository,
\href{https://github.com/sara316129/srilanka-rice-supply-chain-bullwhip}{\textit{Bullwhip Effect
in Sri Lanka's Rice Supply Chain}} on GitHub. See the repository README for reproduction
instructions.

\section*{Acknowledgments}
We thank Dr.\ S.\ C.\ Surasinghe for supervision and guidance throughout this project.

\bibliographystyle{siamplain}
\bibliography{references}

\end{document}